% Natural Language-Driven Constrained Ship Design Decision Framework
% Ocean Engineering manuscript outline based on the Elsevier CAS double-column template.
% Evidence placeholders are intentionally retained where the current project record does
% not yet contain the corresponding experiment.

\documentclass[a4paper,fleqn]{cas-dc}

\usepackage[authoryear,longnamesfirst]{natbib}
\usepackage{amsmath}
\usepackage{booktabs}
\usepackage{graphicx}
\usepackage{xcolor}

\newcommand{\todo}[1]{\textcolor{red}{[TODO: #1]}}
\newcommand{\evidence}[1]{\textcolor{blue}{[Evidence needed: #1]}}

\begin{document}
\let\WriteBookmarks\relax
\def\floatpagepagefraction{1}
\def\textpagefraction{.001}

\shorttitle{Natural-language-driven constrained ship design}
\shortauthors{Author et~al.}

\title [mode = title]{Natural Language-Driven Constrained Ship Design Decision Framework}

% Replace author and affiliation placeholders before submission.
\author[1]{Author Name}
\cormark[1]
\ead{author@example.com}
\affiliation[1]{organization={Institution Name},
                addressline={Address},
                city={City},
                postcode={Postcode},
                country={Country}}
\cortext[1]{Corresponding author}

\begin{abstract}
\todo{Write the structured problem--method--evidence--implication abstract after the
main results and end-to-end evaluation are complete.}
\end{abstract}

\begin{highlights}
\item \todo{State the NURBS-based constrained ship FFD decision problem.}
\item \todo{State the typed natural-language decision framework and action decoder.}
\item \todo{Report the strongest benchmark and geometric validation evidence.}
\end{highlights}

\begin{keywords}
NURBS \sep free-form deformation \sep ship design \sep natural language decision \sep constrained geometric editing
\end{keywords}

\maketitle

% -----------------------------------------------------------------------------
\section{Introduction}
\label{sec:introduction}

Ship-form design requires repeated geometric adjustments under coupled geometric
and engineering constraints. Designers must modify local regions such as the bow,
stern, bulb, deck, or bilge while preserving a valid hull representation and, when
required, global properties such as symmetry, displacement, or the deck line.
NURBS curves provide a compact representation for smooth ship-form geometry, and
free-form deformation (FFD) provides a mechanism for applying localized or global
changes. The remaining design difficulty is to express these geometric operations
as reliable, interpretable decisions that can be selected from natural-language
intent.

Natural-language interfaces are attractive for iterative ship design because a
designer can describe a target change without manually editing control points.
However, unconstrained text generation does not directly guarantee a valid FFD
action: the output must identify a compatible region and operation, select the
number of actions, preserve the semantics of a multi-step interaction, and satisfy
the geometric constraints imposed by the hull representation. This creates a
decision problem at the interface between language understanding and geometric
execution. A model that produces fluent text but selects an invalid action can
still generate an unusable design operation.

This paper formulates that interface as a constrained decision problem for NURBS-
based ship-form FFD. Given a hull state and a natural-language design intent, the
framework represents the candidate operation through typed decisions over action
count, region, operation, and related attributes. A decision model predicts a
probability distribution for each typed question. The maximum-probability answers
are projected into a canonical FFD action representation and passed to a geometric
decoder. The decoder updates NURBS control parameters, reconstructs the hull, and
checks the resulting geometry. This separation keeps language-level decision making
and numerical geometry execution explicit and independently testable.

The study makes the following contributions:
\begin{enumerate}
    \item It defines a NURBS-based ship-form FFD decision problem in which ship
    editing intent is mapped to typed, engineering-interpretable deformation actions.
    \item It develops a natural-language-driven decision framework that combines
    typed probability prediction, action projection, sequential editing, and
    constraint-aware geometric execution.
    \item It constructs a ship design decision dataset and a frozen benchmark with a
    common probability protocol, question-level metrics, calibration metrics, and
    a shared action-level FFD score.
    \item It compares specialized decision models, the native JEV decision model,
    and general language models under the same benchmark protocol. The current
    benchmark results provide the comparison in Section~\ref{sec:benchmark};
    geometric end-to-end and ablation claims remain subject to the experiments
    marked in Sections~\ref{sec:end-to-end} and~\ref{sec:ablations}.
\end{enumerate}

% -----------------------------------------------------------------------------
\section{Methods}
\label{sec:methods}

This section defines the geometric representation, the decision interface, and the
dataset and evaluation protocol. The central design choice is to use the model as a
typed decision layer rather than asking it to emit unrestricted NURBS control-point
coordinates.

\subsection{NURBS-Based Ship Representation and FFD}
\label{sec:nurbs-ffd}

The ship-form representation combines longitudinal waterline curves with fore and
aft contours extracted near the longitudinal centre plane. Non-uniform rational
B-splines (NURBS) describe these curves through control points, positive weights,
and knot vectors. Free-form deformation (FFD) is implemented by updating the
waterline control points, followed by curve evaluation and reconstruction of a
closed hull mesh. This representation connects local geometric edits to a common
set of curve parameters.

\paragraph{Coordinate system and waterline representation.}
The input is a triangulated hull surface. All coordinates are scaled by the hull
length, with the longitudinal origin at the aft extremity, the transverse origin
on the centreline, and the vertical origin at the lowest point of the input hull.
The normalized longitudinal extent is therefore $[0,1]$. Horizontal sections are
extracted at prescribed elevations, and their longitudinal positions and
half-breadths define the waterlines. Each half-waterline is divided at the midpoint
of its longitudinal extent into aft and fore segments. Each segment is fitted by
a cubic, clamped NURBS curve with eight control points:
\begin{equation}
 \mathbf{c}(u)=
 \frac{\sum_{i=0}^{7} B_{i,3}(u)w_i\mathbf{p}_i}
      {\sum_{i=0}^{7} B_{i,3}(u)w_i},
 \qquad 0\leq u\leq 1,
 \label{eq:waterline-nurbs}
\end{equation}
where $\mathbf{p}_i$ is a control point, $w_i>0$ is its weight, and $B_{i,3}$ is
the cubic B-spline basis function defined by the knot vector. The knot vector is
clamped at both ends and has uniformly spaced interior knots.

Both segments are parameterized from the outer hull extremity towards the section
midpoint. Their endpoint positions are fixed during fitting, and three interior
control points coincide to parameterize the transition towards the central
portion of the waterline. Four interior control-point positions and three positive
weights are optimized by bounded nonlinear least squares. The fitting parameters
are obtained from cumulative chord length. After fitting, longitudinal
control-point coordinates are ordered in the direction of the curve. The two
segments are joined at their common midpoint, and the half-waterline is reflected
across the centreline to construct the complete section.

\paragraph{Longitudinal profile representation.}
A plane parallel and close to the longitudinal centre plane intersects the input
mesh to obtain ordered contours in the longitudinal--vertical plane. Contours in
the fore region, with normalized longitudinal coordinate at least $0.65$, and the
aft region, with coordinate at most $0.35$, are assembled and fitted separately.
Each profile uses a cubic, clamped curve with 32 control points and unit weights.
The contour is resampled by arc length, its endpoints are fixed, and the remaining
control-point coordinates are fitted by bounded least squares. The deck and lower
hull portions within each extracted contour are fitted together.

The fore profile incorporates two geometrically identified features in its knot
vector. The most forward position within the lower half of the contour defines a
bulb-related parameter. An upper-contour corner is identified from changes in the
sampled tangent direction. Eight interior knots are concentrated around the
bulb-related parameter, while a knot of multiplicity three is placed at the upper
corner. The remaining interior knots are distributed over the intervening
parameter intervals. The aft profile uses uniformly spaced interior knots. At
each waterline elevation, intersections with the fitted profiles determine the
longitudinal positions of the fore and aft centreline endpoints. The foremost
fore-profile intersection and the aftmost aft-profile intersection are selected
when several intersections occur. These positions couple the longitudinal
profiles to the initial waterline representation.

\paragraph{Control-point deformation.}
A geometric editing action specifies a hull region, an operation, a magnitude,
and optional longitudinal and vertical extents. Supported regions are the bow,
stern, bulb, midbody, deck, bilge, and the complete hull. The magnitude is expressed
as a fraction of the relevant hull dimension; ordinal inputs are mapped to the
fractions $0$, $0.01$, $0.02$, $0.04$, $0.08$, and $0.12$. Before deformation,
the waterline controls are embedded in three dimensions by assigning each control
its section elevation.

Local edits are weighted by the control-point positions within the longitudinal
and vertical action extents. Let $\xi$ and $\zeta$ denote the longitudinal and
vertical coordinates normalized by the current control-point extents. Smooth
transitions are constructed using
\begin{equation}
 S(t)=3\tau^2-2\tau^3,
 \qquad \tau=\min(1,\max(0,t)).
 \label{eq:smooth-transition}
\end{equation}
For a local interval $[a,b]$, the weighting function combines transitions from
both ends:
\begin{equation}
 W(t;a,b)=S\!\left(\frac{t-a}{d}\right)
           S\!\left(\frac{b-t}{d}\right),
 \qquad d=0.2(b-a).
 \label{eq:local-window}
\end{equation}
A full-coordinate interval receives unit weight. At the uppermost and lowermost
vertical limits, the weighting includes the terminal sections when they belong
to the selected region. Multiplying the longitudinal and vertical weights yields
the local influence factor $M$ for each control point.

Transverse expansion and contraction scale the control-point half-breadth by
$1+\alpha M$ and $1-\alpha M$, respectively, where $\alpha$ is the magnitude
fraction. The fullness and flare-increase operations use the same transverse
expansion rule over their selected extents. Vertical edits translate the selected
controls by $\pm\alpha D M$, where $D$ is the current vertical extent.
Longitudinal edits translate controls by a fraction of the current longitudinal
extent. In the bow and stern, a one-sided smooth transition increases the
longitudinal influence towards the corresponding hull extremity; other local
regions use the two-sided weighting function. Global length changes scale
longitudinal coordinates about the centre of the current longitudinal extent,
while global breadth changes scale transverse coordinates. Successive actions
are applied in their specified order to the updated waterline controls.

\paragraph{Constraint handling and hull reconstruction.}
The geometry executor maintains bilateral symmetry by reflecting the deformed
half-waterlines. A deck-line preservation request attenuates local deformation
with increasing normalized elevation and fixes the uppermost controls. When
volume preservation is requested, the pre-edit and post-edit enclosed volumes
are evaluated from the reconstructed meshes, and transverse control-point
coordinates are rescaled by their ratio. With concurrent deck-line preservation,
the transverse compensation is applied below the uppermost controls. Vertical
edits are rejected when fixed draught is requested. An optional shape-change
threshold is evaluated using changes in normalized second differences of the
longitudinal and transverse control-point coordinates.

After each action, the evaluated waterline is checked for longitudinal ordering
and non-negative half-breadth, and its area and centroid are recalculated.
Longitudinal translations include a control-point ordering correction before
curve evaluation. To reconstruct the hull, the aft and fore curves are sampled
and joined, then mirrored to form closed section rings with shared centreline
vertices. Corresponding vertices on adjacent rings are connected by triangles;
the lowest and highest rings are triangulated to close the mesh. The resulting
hull is checked for watertightness and valid enclosed volume. This completes the
editing sequence from NURBS control-point updates to a reconstructed ship form.

\subsection{Decision Framework for Natural-Language-Driven Ship Design}
\label{sec:decision-framework}

\paragraph{Problem interface.}
Define the input as a current hull state together with a natural-language design
intent, and the output as one or more typed deformation actions. Explain how the
intent is represented in the benchmark request and how the state is updated after
sequential turns.

\paragraph{Typed decision layer.}
Describe the three question families used by the framework: \textit{choice} for
categorical selections, \textit{score} for ordinal levels, and \textit{noul} for
Boolean constraints. For a record with questions $q_1,\ldots,q_K$, the model
outputs distributions
\[
 p(y_k\mid s,q_k), \qquad k=1,\ldots,K,
\]
where $s$ is the current state. Explain that the benchmark retains the complete
probability distributions for accuracy, negative log-likelihood, expected
calibration error, and Brier score.

\paragraph{Action projection and execution.}
The argmax answer for each typed question is projected into the canonical action
format
\begin{verbatim}
{"actions":[{"region":"bow","operation":"outward"}]}
\end{verbatim}
The projected actions are then passed to the NURBS/FFD executor. Invalid
probabilities, missing answers, incompatible action counts, failed calls, and
invalid action projections are recorded as failures and included in the relevant
denominators. This section should include a framework diagram showing the flow
from natural-language intent to typed decisions, action projection, NURBS update,
and geometric validation.

\subsection{Ship Design Decision Dataset and Benchmark}
\label{sec:dataset-benchmark}

\paragraph{Dataset construction.}
Describe the programmatic generation of hull states and FFD actions, including
single-action edits, multi-region edits, repeated same-region turns, explicit-range
edits, natural-language paraphrases, and constraint annotations. State the split
policy by hull family and the frozen test-set identifier and checksum.

\paragraph{Benchmark protocol.}
All compared models use the same frozen $5{,}000$-record test set, typed questions,
probability JSON format, probability validation, and metric aggregation. The
frozen set contains 43,496 expanded questions. Report question-level clean
accuracy, negative log-likelihood, expected calibration error, and Brier score;
report FFD action exact-match accuracy separately. Define the turn-level denominator
as the number of requested records, including failed or invalid records.

\paragraph{Compared systems and reproducibility.}
List the Kev-0.8B, Kev-2B, Kev-4B, native JEV, DeepSeek-Flash, Claude Opus 5,
GPT-5.6 Sol, Qwen3, Gemma3, and Llama3.2 systems. State model identifiers,
provider or local-serving mode, decoding/thinking setting, checkpoint information
for local specialized models, and code/data revisions in the final manuscript.

% -----------------------------------------------------------------------------
\section{Results}
\label{sec:results}

\subsection{NURBS Representation and FFD Accuracy}
\label{sec:nurbs-results}

\paragraph{Claim to test.}
The NURBS representation should reconstruct valid ship forms with sufficiently low
geometric error, while the FFD executor should produce valid deformations under the
specified action and constraint conditions.

\paragraph{Planned evidence.}
Report fitting and reconstruction errors, watertightness and volume validity, and
representative before/after hull visualizations. Include a table separating
representation error from post-FFD validity so that geometric fidelity is not
confounded with decision accuracy.

\evidence{Final quantitative reconstruction and FFD execution results.}

\subsection{Ship Design Decision Benchmark Evaluation}
\label{sec:benchmark}

\paragraph{Claim supported by current results.}
Under the shared frozen benchmark, specialized Kev models and several language
models can be compared directly through the same question-level probability metrics;
FFD action exact-match is reported as a separate projection metric.

\begin{table*}[t]
\caption{Outline of the shared frozen test-set comparison. Values currently
available from the local reports are shown; negative log-likelihood, expected
calibration error, and Brier score should be added from the final report extraction.}
\label{tab:benchmark-outline}
\centering
\begin{tabular}{lrrrr}
\toprule
Model & Clean accuracy & NLL & ECE & FFD exact match \\
\midrule
Kev-0.8B       & 95.90\% & \todo{add} & \todo{add} & 99.32\% \\
Kev-2B         & 95.92\% & \todo{add} & \todo{add} & 99.30\% \\
Kev-4B         & 95.88\% & \todo{add} & \todo{add} & 99.24\% \\
JEV            & 90.96\% & \todo{add} & \todo{add} & 84.52\% \\
DeepSeek-Flash & 92.84\% & \todo{add} & \todo{add} & 98.40\% \\
Claude Opus 5  & 92.87\% & \todo{add} & \todo{add} & 99.12\% \\
GPT-5.6 Sol    & 92.44\% & \todo{add} & \todo{add} & 98.32\% \\
Qwen3          & 75.61\% & \todo{add} & \todo{add} & 2.84\% \\
Gemma3         & 66.97\% & \todo{add} & \todo{add} & 15.98\% \\
Llama3.2       & 63.94\% & \todo{add} & \todo{add} & 2.16\% \\
\bottomrule
\end{tabular}
\end{table*}

Discuss coverage and failure rates alongside accuracy. The table must distinguish
question-level clean accuracy from turn-level FFD exact match; these metrics have
different statistical objects and must not be merged into one ranking.

\subsection{End-to-End Ship Form Editing Evaluation}
\label{sec:end-to-end}

\paragraph{Claim to test.}
A correct typed decision should produce a valid and interpretable geometric edit
when decoded by the NURBS/FFD executor.

\paragraph{Planned protocol.}
Evaluate complete natural-language instruction to edited-hull trajectories,
including single-turn, multi-region, repeated-turn, and reversal cases. Report
instruction success, geometric validity, constraint satisfaction, displacement or
principal-dimension changes, and human or expert agreement where available.

\evidence{End-to-end execution runs, geometric metrics, representative hull
figures, and the definition of the final acceptance criteria.}

\subsection{Ablation Studies}
\label{sec:ablations}

Each ablation should use the same frozen split and failure denominator as the main
benchmark. The reported change should be tied to one design choice at a time.

\subsubsection{Sensitivity to NURBS Representation}
\label{sec:ablation-nurbs}
Compare control-point count, spline degree, waterline sampling, and representation
variants while keeping the decision protocol fixed. Report both reconstruction
quality and downstream action-to-geometry validity.

\evidence{Representation variants and their measured effect.}

\subsubsection{Effect of Kev Model Scale}
\label{sec:ablation-scale}
Compare Kev-0.8B, Kev-2B, and Kev-4B under identical data, prompts, evaluation
code, and test records. Report clean accuracy, calibration, FFD exact match,
coverage, and latency.

\evidence{Final scale-comparison analysis, including uncertainty or repeated-run
statistics if available.}

\subsubsection{Effect of Training Data Size}
\label{sec:ablation-data}
Evaluate checkpoints trained with controlled fractions of the training data. Keep
architecture and optimization settings fixed, and plot performance against the
number of training examples.

\evidence{Controlled data-size checkpoints and a reproducible training-data
schedule.}

\subsubsection{Generalization and Robustness}
\label{sec:ablation-robustness}
Test unseen hull families, paraphrased instructions, multi-step edits, constraint
variants, and malformed or ambiguous requests. Report both successful decisions
and failure categories rather than aggregating all failures into a single number.

\evidence{Held-out hull-family, paraphrase, multi-step, and stress-test results.}

% -----------------------------------------------------------------------------
\section{Discussion}
\label{sec:discussion}

\subsection{Why LLMs Fail in Ship FFD Decision Making}
\label{sec:llm-failures}
Interpret failure categories revealed by the benchmark, including action-count
errors, region/operation confusions, invalid probability structures, unsupported
multi-action projections, and failures to preserve the typed semantics of the
request. Connect each explanation to error statistics or qualitative examples.

\subsection{Specialized Model Scale and Data Efficiency}
\label{sec:scale-efficiency}
Discuss the relationship between model scale, training data, calibration, and
action-level reliability. Separate observations directly supported by the scale
comparison from hypotheses about why specialized decision training helps.

% -----------------------------------------------------------------------------
\section{Conclusion and Future Work}
\label{sec:conclusion}

Summarize the bounded contribution: a NURBS-based constrained ship FFD decision
problem and a natural-language-driven typed decision framework that can be
assessed with a common probability and action protocol. State the strongest
validated benchmark evidence, then identify the current boundary: the shared
benchmark primarily evaluates typed decisions and their action projection, while
full continuous parameter control, hydrodynamic objectives, and validated
end-to-end design optimization require further experiments.

Future work should address explicit continuous magnitudes and spatial extents,
stronger constraint and hydrodynamic coupling, expert-reviewed multi-turn design
trajectories, unseen-hull generalization, and calibrated human-in-the-loop
acceptance or rollback.

% -----------------------------------------------------------------------------
% Add references after the literature review has been completed.
% \bibliographystyle{cas-model2-names}
% \bibliography{references}

\end{document}
